\documentstyle[% 


righttag% 


,11pt% 


,amssymb% 


,russian%               remove in Eng. 


,izvru%                 Set izven% in Eng. 


]{amsart} 


 


%       Math definitions 


 


\newcommand{\const}{\operatorname{const}} 


\newcommand{\sign}{\operatorname{sign}} 


\newcommand{\tts}[1]{} 


 


\theoremstyle{plain} 


\theorembodyfont{\it} 


\newtheorem{thmnonum}{\hskip\parindent Теорема} 


\renewcommand{\thethmnonum}{} 


 


\theoremstyle{definition} 


\theorembodyfont{\rm} 


\newtheorem{notenonum}{\hskip\parindent Замечание} 


\renewcommand{\thenotenonum}{} 


\renewcommand{\baselinestretch}{0.97}


 


%\hoffset=-15mm%                  For Eng. 


  


\setcounter{page}{44} 


\god{1999} 


\nomer{4~(443)} %                        Remove in Eng. 


%\nomer{Vol.\,43, No.\,4}%               Set in Eng. 


%\str{\pageref{begin}--\pageref{end}}%  Set in Eng. 


\udk{517.927} 


 


\author{\it А.Л.\,Тептин} 


% NB:   ^^ remove in Eng. 


\title{К вопросу об осцилляционности спектра\\ 


многоточечной краевой задачи} 


 


\date{} 


% \pagestyle{myheadings}%         Set in Eng. 


\pagestyle{plain} %              Remove in Eng. 


\begin{document} 


% \thispagestyle{plain}%          Set in Eng. 


\maketitle 


\label{begin} 


 


 


Рассмотрим задачу на собственные значения 


\begin{gather} 


L y\equiv y^{(n)}+\sum^{n-1}_{k=0}p_k(\cdot)y^{(k)}= 


\lambda q(\cdot)y, \\ 


% 


l_{j c_j}y\equiv\sum^n_{k=1}b_{jk}y^{(n-k)}(c_j)=0, 


\quad j=\ov{1,r}, \\ 


% 


y^{(\nu_i)}(a_i)=0,\quad \nu_i=\ov{0,k_i-1},\quad 


i=\ov{1,m},


\end{gather} 


где $1<r\le n$, $1\le m\le n-r$, $k_1+\dotsb+k_m=n-r$, $b_{jk}\in\cal 


R$, $k=\ov{1,n}$, 


$b_{j1}=1$, $j=\ov{1,r}$, $a\le c_1\le\dotsb\le c_r\le b$, $a\le 


a_1<\dotsb<a_m\le b$, 


\begin{equation} 


\min\{c_1,a_1\}=a,\quad \max\{c_r,a_m\}=b. 


\end{equation} 


 


Следуя [1], спектр задачи (1)--(3) будем называть осцилляционным, 


если он состоит из собственных значений, положительных простых и 


образующих неограниченную последовательность 


$0<\lambda_0<\lambda_1<\dotsb$, причем 


для каждого $k$, $k=0,1,\dotsc$, собственная функция $y_k(x)$ задачи 


(1)--(3), 


соответствующая собственному значению $\lambda_k$, такова, что функция 


$\psi_k(x)=y_k(x)/w(x)$, 


где $w(x)=\prod\limits^m_{i=1}(x-a_i)^{k_i}$ при $r<n$ и $w(x)\equiv 1$ 


при $r=n$, 


непрерывна в $[a,b]$, имеет в $[a,b]$ ровно $k$ нулей, которые все 


изолированные, узловые, нули $\psi_k(x)$ и $\psi_{k+1}(x)$ 


перемежаются, а последовательность 


$\{\psi_k(x)\}^\infty_0$ образует в $[a,b]$ ряд Маркова ([2]--[4]). 


 


Условия осцилляционности спектра задачи (1)--(3) при $r=0$ рассмотрены, 


например, в ([1], [5], [6]), при $r=1$ --- в [7], при совпадении каждой 


из точек $c_j$ либо с $a$, либо с $b$ --- в [3], [8]--[10], 


при $b_{jk}=0$, $k=\ov{1,j-1}$, $j=\ov{1,r}$ --- в [6], [11], 


при $b_{j1}\ne 0$, $j=\ov{1,r}$, произвольном $r$, $0<r<n$, $m\ge2$, 


$a=a_1<a_2<\dotsb<a_m=b$, $a<c_1<c_2<\dotsb<c_r<b$ 


либо $a=c_1\le c_2\le\dotsb\le c_r=b$ --- в [12]--[14]. 


 


В данной статье результаты работ [12]--[14] распространяются 


на любые случаи расположения точек $a_i$, $i=\ov{1,m}$, $c_j$, 


$j=\ov{1,r}$, в $[a,b]$ при 


соблюдении условия (4), любых $r$, $0<r\le n$, и $m$, $m\ge 1$, а также 


обсуждаются 


свойства, которыми будет обладать спектр задачи (1)--(3), 


если для миноров матрицы $\cal B=(b_{jk})$, $j=\ov{1,r}$, $k=\ov{1,m}$, 


требуемые в [10], [12]--[14] 


неравенства 


\begin{equation} 


\cal B\binom{\nu\dotsb\nu+\mu}{1\dotsb\mu+1}>0,\quad 


\nu=\ov{1,r-\mu},\ \; \mu=\ov{1,r-1}, 


\end{equation} 


гарантирующие при прочих условиях этих работ для функции Грина 


$\cal G(x,s)$ уравнения 


\begin{equation} 


L y=0 


\end{equation} 


с краевыми условиями (2), (3) знакопостоянство функции 


$\cal G(x,s)w(x)$ 


[15] в каждой из непустых полос 


\begin{equation} 


a\le x\le b,\quad c_j<s<c_{j+1},\quad 


j=\ov{0,r},\quad c_0=a,\quad c_{r+1}=b, 


\end{equation} 


заменить существованием перестановки $(j_1,j_2,\dotsc,j_r)$ множества 


$\{1,2,\dotsc,r\}$, 


для которой 


\begin{gather} 


j_\nu<j_{\nu+k},\quad k=\ov{2,r-\nu},\quad \nu=\ov{1,r-2}, \\ 


\cal B\binom{j_\nu\dotsb j_{\nu+\mu}}{1\dotsb\mu+1}>0,\quad 


\nu=\ov{1,r-\mu},\quad \mu=\ov{1,r-1}, 


\end{gather} 


что при сохранении прочих условий работ [10], [12]--[14] также 


обеспечивает 


сохранение знака $\cal G(x,s)w(x)$ в каждой из полос (7) [16], но, 


возможно, 


иного, нежели при условии (5). 


\medskip 


 


{\bf 1.} Как известно, нулевым местом или отдельным нулем непрерывной 


функции $u(x)$ называется всякая компонента связности множества 


ее нулевых точек. Нулевое место $[c,d]\subset(a,b)$ называется узлом 


(пучностью), если $u(c-\varepsilon)u(d+\varepsilon)<0$ ($>0$) для всех 


достаточно малых 


$\varepsilon>0$ ([2]--[4]). Следуя [1], $k$-кратностью изолированного 


нуля $x_0$ 


функции $u(x)$ назовем число $r_k(u,x_0)=\max\{i\le 


k:u(x)=o(|x-x_0|^{i-1})\}$, 


а для нулевого места $\Omega$ другого типа положим $r_k(u,\Omega)=k$. 


Пусть $r_k(u,\cal J)$ --- сумма $k$-кратностей нулевых мест 


$u(x)$ в промежутке $\cal J$, $r_k(u)=r_k(u,[a,b])$, так что $r_1(u)$ --- 


число отдельных нулей $u(x)$ в $[a,b]$; $\varphi(u,\cal J)$ --- число 


нулей $u(x)$ на 


$\cal J$ с учетом их обычной кратности [17], 


$\varphi(u)=\varphi(u,[a,b])$; 


$S(u,\cal J)$ --- число перемен знака $u(x)$ на $\cal J$, т.\,е.~такое 


число $l$, 


что $\exists C=\const\ne 0$, $\exists x_i\in \cal J$, $i=\ov{1,l}$, 


$\inf\cal J=x_0<x_1<\dotsb<x_{l+1}=\sup\cal J:(-1)^i C u(x)\ge 0$ 


почти всюду в $(x_i,x_{i+1})$, $u(x)\nsim 0$ в $(x_i,x_{i+1})$, 


$i=\ov{0,l}$, при этом 


в случае $\cal J=[a,b]$ положим $\sign_1u=\sign C$ ([3], [4], [8]), 


$\sign_2u=(-1)^l\sign C$, $S(u)=S(u,[a,b])$ [3], [4]; $C^k\cal J$ ($\wt 


C^k\cal J$) --- 


множество вещественных функций, имеющих в $\cal J$ непрерывную 


(абсолютно 


непрерывную) $k$-ю производную, $\wt C^{-1}\cal J$ --- множество 


суммируемых в $\cal J$ функций; 


$C_*[a,b]$ --- множество непрерывных на $[a,b]$ функций, обладающих 


только изолированными нулями; $\delta[a,b]$ --- множество обобщенных 


функций вида $u(x)=\sum\limits^k_{i=1}A_i\delta(x-s_i)$, 


$a<s_1<\dotsb<s_k<b$, $k=1,2,\dotsc$ ($\delta(x)$ --- дельта-функция), 


для которых будем полагать $S(u)=S(U)$, 


$\sign_1u=\sign_1U$, $\sign_2u=\sign_2U$, где $U=(A_1,\dotsc,A_k)$ [3], 


[4], а $S(U)$ 


для ненулевого вектора $U$ будет означать число перемен знака в ряду 


его координат с выброшенными нулями, при этом $\sign_1U=\sign{}A^*$, 


$\sign_2U=\sign{}A^{**}$, где $A^*$ --- первая, а $A^{**}$ 


--- последняя ненулевые


координаты $U$ [3], [4]; $\cal B(j_1\dotsb j_\mu)=\cal 


B\binom{j_1\dotsb j_\mu}{1\dotsb\mu}$ --- минор матрицы $\cal B$; 


$\cal D(j_1\dotsb j_{\mu+1})=\cal B(j_1\dotsb j_{\mu-1})\cal B(j_1\dotsb 


j_{\mu+1})/(\cal B(j_1\dotsb j_\mu)\cal B(j_1\dotsb 


j_{\mu-1}j_{\mu+1}))$, $\mu=\ov{2,r-1}$, $\cal D(j_1j_2)=\cal 


B(j_1j_2)$; $c=(c_1,\dotsc,c_r)$ --- 


вектор, $c_0=a$, $c_{r+1}=b$; $\sigma_c(s)=(-1)^{r-j}$ 


для $s\in(c_j,c_{j+1})$, $j=\ov{0,r}$, $\sigma_c(c_j)$, $j=\ov{1,r}$, 


будем 


определять произвольно; $I$ --- единичный оператор; $\cal L^p_k$ --- 


множество 


дифференциальных операторов вида 


$\frac{d^k}{dx^k}+\sum\limits^{k-1}_{i=0} 


p_i(\cdot)\frac{d^i}{dx^i}$, $p_i(\cdot)\in\wt C^{p-1}[a,b]$, 


$i=\ov{0,k-1}$, $\cal L^p_0=\{I\}$; $T^p_k(a,b)$ --- 


класс таких операторов $\cal V\in\cal L^p_k$, что $\varphi(y)\le k-1$ 


для любого решения $y(x)\not\equiv0$ уравнения $\cal V y=0$, 


$T^p_0(a,b)=\cal L^p_0$; $\wt T^p_n(\cal B;a,b;c_1,\dotsc,c_r)$ при 


данных $\cal B$ и $c_j$, $j=\ov{1,r}$, --- 


класс таких операторов $L\in\cal L^p_n$, что $\varphi(y)\le n-\mu-1$ 


для любого решения $y(x)\not\equiv0$ 


уравнения (6), удовлетворяющего любым $\mu$ из условий 


(2) $\forall\mu$, $\mu=\ov{1,r}$; $T^p_{n1}(\cal B;a,b;c_1,c_r)$ --- 


пересечение классов 


$\wt T^p_n(\cal B;a,b;c'_1,\dotsc,c'_r)$ для всевозможных $c'_j$, 


$j=\ov{1,r}$, $c_1\le c'_1\le\dotsb\le c'_r\le c_r$; 


$Q=[a,b]\times([a,b]\setminus\{c_j\}^r_1)$. 


 


Очевидно, для данного вектора $c=(c_1,\dotsc,c_r)$ при условии 


$L\in\wt T^0_n(\cal B;a,b;c_1,\dotsc,c_r)$ существует функция Грина 


$\cal G_c(x,s)$ 


задачи (6), (2), (3). Пусть 


\begin{equation} 


\Gamma_c(x,s)= 


\begin{cases} 


\cal G_c(x,s)\sigma_c(s)/w(x) & \text{при \ }x\ne a_i, \\ 


(\partial^{k_i}\cal G_c(a_i,s)/\partial x^{k_i}) 


\sigma_c(s)/w^{(k_i)}(a_i)& \text{при \ }x=a_i,\ \; i=\ov{1,m}. 


\end{cases} 


\end{equation} 


Непрерывность $\Gamma_c(x,s)$ в $Q$ 


и существование равномерных конечных пределов $\Gamma_c(x,c_j\pm 0)$, 


$j=\ov{1,r}$, можно доказать, дословно 


повторяя рассуждения 


из [13], [14], которые применимы здесь, т.\,к.~$k_i\le n-r\le n-2$ 


при $r>1$, и справедливы при любом расположении 


точек $a_i$, $i=\ov{1,m}$, и $c_j$, $j=\ov{1,r}$, на $[a,b]$. Поэтому 


$\Gamma_c(x,\cdot)$ и $\cal G_c(x,\cdot)$ 


при $s=c_j$, $j=\ov{0,r+1}$, можно доопределить по непрерывности с 


любой из сторон либо следующим образом: 


$\Gamma_c(x,c_j)=0.5(\Gamma_c(x,c_j+0)+ 


\Gamma_c(x,c_j-0))$, $\cal G_c(x,c_j)=0.5(\cal G_c(x,c_j+0)+\cal 


G_c(x,c_j-0)\sigma_c(c_j-0)\sigma_c(c_j+0))$, $a\le x\le b$, 


$c_j\in(a,b)$, $1\le j\le r$. 


 


В силу (10), теоремы [15], леммы 1 [13] (справедливость которой 


не зависит от строгости или нестрогости неравенств $c_j\le c_{j+1}$, 


$j=\ov{0,r}$) 


и замечания к ней из [14] при условии (5) и 


\begin{equation} 


L\in T^0_{n1}(\cal B;a,b;c_1,c_r) 


\end{equation} 


выполняется неравенство 


\begin{equation} 


\Gamma_c(x,s)>0\quad 


\text{в} 


\quad [a,b]\times(a,b). 


\end{equation} 


 


{\bf 2.} В силу теорем 1 и 2 [18] при условии (11) для каждого вектора 


$c=(c_1,\dotsc,c_r)$, $a\le c_1\le\dotsb\le c_r\le b$ и каждой 


перестановки $(j_1,\dotsc,j_r)$ 


множества $\{1,\dotsc,r\}$ оператор $L$ допускает в $[a,b]$ разложение 


\begin{equation} 


L^{j_1\dotsb j_\mu c}_{n-\mu}=\bigg( 


\frac{d}{dx}-\alpha^c_{j_1\dotsb j_{\mu+1}}(x)I\bigg) 


L^{j_1\dotsb j_{\mu+1}c}_{n-\mu-1},\quad \mu=\ov{0,r-1}, 


\end{equation} 


где $L^{j_1\dotsb j_\mu c}_{n-\mu}=L$ при $\mu=0$, 


$L^{j_1\dotsb j_\mu c}_{n-\mu}$ не зависят от 


$c_{j_{\mu+1}},\dotsc,c_{j_r}$ 


и от 


порядка индексов $j_1,\dotsc,j_\mu$, $\mu=\ov{1,r}$, 


\begin{gather} 


L^{j_1\dotsb j_r c}_{n-r}=L^c_{n-r}\in 


T^r_{n-r}(a,b),\\ 


(L^{jc}_{n-1}y)(c_j)=l_{j c_j}y\ \; \forall y\in C^{n-1} 


[a,b],\quad j=\ov{1,r}, 


\end{gather} 


$\alpha^c_{j_1\dotsb j_\mu}(x)$ 


не зависят от $c_{j_{\mu+1}},\dotsc,c_{j_r}$ и от порядка индексов 


$j_1,\dotsc,j_{\mu-1}$, $\mu=\ov{2,r}$, 


$\alpha^c_{j_1\dotsb j_\mu}(x)\in\wt C^{\mu-2}[a,b]$, 


$\mu=\ov{1,r}$, 


\begin{equation} 


(\alpha^c_{j_1\dotsb j_{\mu+1}}(x)- 


\alpha^c_{j_1\dotsb j_{\mu-1}j_{\mu+1}j_\mu}(x))\cal D 


(j_1,\dotsc,j_{\mu+1})>0 \ \; \forall x\in[a,b],\ \; \mu=\ov{1,r-1}. 


\end{equation} 


Кроме того, из (14) (см., напр., [17]) следует факторизация 


\begin{equation} 


L^c_k=\bigg(\frac{d}{dx}-\beta_k(x)I\bigg) 


L^c_{k-1},\quad \beta_k(\cdot)\in \wt C^{n-k-1}[a,b], 


\quad k=\ov{1,n-r},\quad L^c_0=I. 


\end{equation} 


 


Пусть 


\begin{gather} 


\Gamma_c z=\int^b_a\Gamma_c(\cdot,s)z(s)ds, \\ 


y_c(x)=(\Gamma_c u)(x),\\ 


v_c(x)=y_c(x)w(x), \\ 


v^c_0(x)=v^{0c}_{j_1\dotsb j_r}(x)=(L^c_{n-r}v_c)(x), \\ 


v^{r-\mu c}_{j_1\dotsb j_\mu}(x)=\bigg( 


\frac{d}{dx}-\alpha^c_{j_1\dotsb j_{\mu+1}}(x)I\bigg) 


v^{r-\mu-1c}_{j_1\dotsb j_{\mu+1}}(x)= 


(L^{j_1\dotsb j_\mu c}_{n-\mu}v_c)(x), 


\ \; \mu=\ov{0,r-1}, 


\end{gather} 


так что в силу (10), (13), (18)--(22) при $\mu=0$\; 


$v^{r-\mu c}_{j_1\dotsb j_\mu}(x)=v^{r c}(x)=(L v_c)(x)=\sigma_c(x) 


u(x)$. 


 


В силу (22) и свойств операторов 


$L^{j_1\dotsb j_\mu c}_{n-\mu}$ функции 


$v^{r-\mu c}_{j_1\dotsb j_\mu}(x)$ не 


зависят от порядка индексов $j_1,\dotsc,j_\mu$, $\mu=\ov{1,r}$. 


Поэтому из (22) следует 


$$ 


v^{r-\mu c}_{\nu\dotsb\nu+\mu-1}(x)- 


v^{r-\mu c}_{\nu+1\dotsb\nu+\mu}(x)= 


(\alpha^c_{\nu+1\dotsb\nu+\mu\nu}(x)- 


\alpha^c_{\nu\dotsb\nu+\mu}(x)) 


v^{r-\mu-1c}_{\nu\dotsb\nu+\mu}(x), 


$$ 


где при условиях (5), (11) в силу (16) 


\begin{equation} 


\alpha^c_{\nu+1\dotsb\nu+\mu\nu}(x)- 


\alpha^c_{\nu\dotsb\nu+\mu}(x)<0 


\ \; \forall x\in[a,b],\ \; \nu=\ov{1,r-\mu}, 


\ \; \mu=\ov{1,r-1}. 


\end{equation} 


 


Согласно лемме 3 [14] при условии (11) все функции 


$v^{r-\mu c}_{j_1\dotsb j_\mu}(x)$, $\mu=\ov{1,r}$, 


непрерывны по $(x,c)$ в области $\cal H=\{(x,c_1,\dotsc,c_r): 


a\le x\le b,\ a\le c_1\le\dotsb\le c_r\le b\}$. 


\medskip 


 


{\bf 3. Лемма 1.}{\it  


Если $S(u)=p$, 


\begin{equation} 


u(\cdot)\in C_*[a,b] 


\end{equation} 


и все узлы $u(x)$ отличны от $c_j$, $j=\ov{1,r}$, то при условиях 


$(5)$, $(11)$ либо 


\begin{equation} 


S(v^c_0)\le r_1(v^c_0)\le S(u)=p, 


\end{equation} 


либо 


\begin{equation} 


r_1(v^c_0)=p+1,\ \; v^c_0(a)=0,\ \; v^c_0(b)=0, 


\ \; S(v^c_0)\le r_1(v^c_0,(a,b))<p. 


\end{equation} 


Если при этом 


\begin{equation} 


S(v^c_0)=S(u), 


\end{equation} 


то 


\begin{equation} 


\sign_1v^c_0=\sign_1u. 


\end{equation} 


} 


 


\begin{pf} При $a<c_1<\dotsb<c_r<b$ и $a=c_1\le\dotsb\le c_r=b$ 


справедливость леммы доказана в [12]--[14], причем установлено 


выполнение (25). 


 


Пусть $a\le c_1\le\dotsb\le c_r\le b$, причем выполняется одно из 


неравенств $a<c_1$ или $c_r<b$. 


 


Как в [14], выберем $\varepsilon_0>0$ меньше каждого из следующих чисел: 


1)~$k$-й доли расстояния от каждой из точек, в которой совпадают 


$k$ последовательных значений $c_j$, до ближайшего к ней узла функции 


$u(x)$; 2)~расстояния от $a$ до первого и от $b$ до последнего узла 


каждой из функций $u(x)$ и $v^{r-\mu c}_{\nu\dotsb\nu+\mu-1}(x)$, 


$\nu=\ov{1,r-\mu+1}$, $\mu=\ov{1,r}$; 3)~$\frac{c_{j+1}-c_j}{m_1+m_2}$ 


для всех $j$, $0\le j\le r$, при которых 


$c_{j-m_1+1}=\dotsb=c_j<c_{j+1}=\dotsb=c_{j+m_2}$. 


 


Пусть $c'_j=a+j\varepsilon$, $j=\ov{0,k}$, если $c_j=a$, $j=\ov{0,k}$, 


$c_{k+1}\ne a$; $c'_j=b-(r-j+1)\varepsilon$, $j=\ov{r-k+1,r+1}$, если 


$c_j=b$, $j=\ov{r-k+1,r+1}$, $c_{r-k}\ne b$; 


$c'_{j+i}=c_j-([k/2]-i)\varepsilon$, $i=\ov{0,k-1}$, если 


$c_{j-1}<c_j=\dotsb=c_{j+k-1}<c_{j+k}$, $1\le j\le r-k+1$, где 


$0<\varepsilon<\varepsilon_0$; $c'=(c'_1,\dotsc,c'_r)$. 


 


В силу теоремы 1 [13] или леммы 1 [12] 


$S(v^{r-\mu c'}_{\nu\dotsb\nu+\mu-1})\le p+r-\mu$, $\nu=\ov{1,r-\mu+1}$, 


$\mu=\ov{1,r}$, 


откуда в силу непрерывности $v^{r-\mu c'}_{\nu\dotsb\nu+\mu-1}(x)$ в 


$\cal H$ 


при достаточно малом $\varepsilon$\; 


$S(v^{r-\mu c}_{\nu\dotsb\nu+\mu-1})\le S(v^{r-\mu 


c'}_{\nu\dotsb\nu+\mu-1})\le p+r-\mu$, $\nu=\ov{1,r-\mu+1}$, 


$\mu=\ov{1,r}$, 


так что 


\begin{equation} 


S(v^c_0)\le S(v^{c'}_0)\le p, 


\ \; S(v^{1c}_{1\dotsb r-1})\le S(v^{1c'}_{1\dotsb r-1}) 


\le p+1, \ \; S(v^{1c}_{2\dotsb r})\le 


S(v^{1c'}_{2\dotsb r})\le p+1. 


\end{equation} 


Так как при условии (24) все нули всех 


$v^{r-\mu c}_{j_1\dotsb j_\mu}(x)$, $\mu=\ov{1,r}$, изолированные 


[10], то в силу (29) и леммы 2 [13] $r_1(v^c_0)\le p+2$. 


 


Пусть для фиксированного вектора $c=(c_1,\dotsc,c_r)$ 


\begin{gather} 


\wt L^{j_1c}_{r-1}:=\bigg(\frac{d}{dx}- 


\alpha^c_{j_1j_2}(x)I\bigg)\dotsb\bigg(\frac{d}{dx} 


-\alpha^c_{j_1\dotsb j_r}(x)I\bigg), \\ 


\wt L^c_r:=\bigg(\frac{d}{dx}-\alpha^c_{j_1}(x)I\bigg) 


\wt L^{j_1c}_{r-1},\\ 


\wt l_{j c_j}y:=(\wt L^{j c}_{r-1}y)(c_j),\quad j=\ov{1,r}, \\ 


\cal G_{r c}(\cdot,s)=L^c_{n-r}\cal G_c(\cdot,s),\quad 


\Gamma_{r c}(x,s)=\cal G_{r c}(x,s)\sigma_c(s), 


\end{gather} 


так что в силу (13) $L=\wt L^c_r L^c_{n-r}$. 


 


Для любой $u(\cdot)\in C_*[a,b]$, узлы которой отличны от $c_j$, 


$j=\ov{1,r}$, 


всегда можно построить такую функцию $u_1(\cdot)\in C_*[a,b]$, что 


$|u_1(x)|\le 1$ 


$\forall x\in[a,b]$, $u_1(x)>0$ почти всюду в $[a,b]$ и при 


\begin{equation} 


u_{\pm\delta}(x)=u(x)\pm\delta u_1(x) 


\end{equation} 


будет 


\begin{equation} 


u_{\pm\delta}(\cdot)\in C_*[a,b],\quad 


S(u)=S(u_{\pm\delta}) 


\end{equation} 


для любого достаточно малого $\delta>0$, причем все узлы 


$u_{\pm\delta}(x)$ отличны 


от $c_j$, $j=\ov{1,r}$. 


 


Пусть $v^c_{01}(x)$ --- функция, построенная по формулам (18)--(21) для 


$u(x):=u_1(x)$, непрерывная в $[a,b]$ $\forall c$. Ввиду (10), (33) 


\begin{multline} 


v^c_{01}(x)=L^c_{n-r}\bigg(w(\cdot)\int^b_a 


\Gamma_c(\cdot,s)u_1(s)ds\bigg)=\int^b_a 


L^c_{n-r}\cal G_c(\cdot,s)\sigma_c(s)u_1(s)ds= \\


= \int^b_a\cal G_{r c}(\cdot,s)\sigma_c(s)u_1(s)ds. 


\end{multline} 


 


В [13] показано, что $\cal G_r(x,s)$ есть функция Грина задачи 


$$ 


\wt L^c_r y=0,\quad \wt l_{j c_j}y=0,\quad 


j=\ov{1,r}. 


$$ 


Как видно из доказательства теоремы [15], справедливость ее заключения 


обеспечивается в конечном итоге факторизацией (13)--(15) и неравенствами 


(23). Так как факторизация (30)--(32) того же типа, что и (13)--(15), 


а коэффициенты $\alpha^c_{j_1\dotsb j_\mu}(x)$, $\mu=\ov{1,r}$, в (13) 


и (30), (31) одни и те же, то заключение теоремы [15] верно и для $\cal 


G_{r c}(x,s)$, т.\,е. 


\begin{gather*} 


\cal G_{r c}(x,s)\sigma_c(s)>0, \ \;x\in[a,b], 


\ \; s\in\bigcup^{r-1}_{j=1}(c_j,c_{j+1}),\\ 


\cal G_{r c}(x,s)=0,\ \; a\le x\le s, 


\ \; \cal G_{r c}(x,s)\sigma_c(s)>0,\ \;s<x\le b, 


\ \; c_r<s\le b, \\ 


\cal G_{r c}(x,s)\sigma_c(s)>0,\ \; a\le x<s, 


\ \; \cal G_{r c}(x,s)=0,\ \; s\le x\le b, 


\ \; a\le s<c_1. 


\end{gather*} 


Отсюда в силу (36) 


\begin{equation} 


v^c_{01}(x)>0 \ \; \forall x\in[a,b]. 


\end{equation} 


 


Если для данного $u(x)$ 


\begin{equation} 


r_1(v^c_0)>p, 


\end{equation} 


то в силу (29) $v^c_0(x)$ имеет нуль хотя бы в одной из точек $x=a$, 


$x=b$ 


или (и) имеет в $(a,b)$ хотя бы один нуль-пучность. Но если $x_i$, 


$i=\ov{1,p_0}$, --- 


все узлы $v^{1c_1}_{1\dotsb r-1}(x)$ на $[a,b]$, $x_i<x_{i+1}$, 


$i=\ov{1,p_0-1}$, $x_0=a$, $x_{p_0+1}=b$, то 


ввиду изолированности всех нулей $v^c_0(x)$ и леммы 2 [13] эта функция 


имеет не более одного нуля в каждом из сегментов $[x_i,x_{i+1}]$, 


$i=\ov{0,p_0}$, 


причем нуль в $(x_i,x_{i+1})$, $0\le i\le p_0$, может быть только узлом. 


Значит, если 


$v^c_0(x)$ имеет нуль-пучность, то он совпадает с одной из точек $x_i$ 


$\exists i=i_0$, $1\le i_0\le p_0$, 


поэтому в силу (22) и леммы 2 [13] 


$\exists\cal D:\cal D v^c_0(x)>0$ $\forall 


x\in[x_{i_0-1},x_{i_0+1}]\setminus\{x_{i_0}\}$. 


Таким образом, в силу (29) в случае (38) $v^c_0(x)$ может 


иметь не более одного нуля-пучности, причем при наличии последнего 


\begin{equation} 


r_1(v^c_0)=p+1. 


\end{equation} 


 


Если 


\begin{equation} 


r_1(v^c_0)=p+2, 


\end{equation} 


то $v^c_0(x)$ не имеет нуля-пучности в $[a,b]$ и при достаточно малом 


$\delta>0$ 


хотя бы для одной из функций $v^c_{0,\pm\delta}(x)$, построенной по 


формулам (18)--(21)


для $u(x):=u_{\pm\delta}(x)$ (см. (34)), т.\,е.~для 


\begin{equation} 


v^c_{0,\pm\delta}(x)=v^c_0(x)\pm\delta 


v^c_{01}(x) 


\end{equation} 


в силу (35), (37) и непрерывности $v^c_0(x)$ и $v^c_{01}(x)$ на 


$[a,b]$\; $S(v^c_{0,\pm\delta})\ge r(v^c_0)-1=p+1=S(u_{\pm\delta})+1$, 


что противоречит (29). Значит, случай 


(40) невозможен. 


 


Аналогичным образом при достаточно малом $\delta>0$ хотя бы для одной 


из функций $v^c_{0,\pm\delta}(x)$ в случае (39) при $v^c_0(a)\ne0$ или 


$v^c_0(b)\ne0$, в том 


числе и при наличии нуля-пучности $v^c_0(x)$ в $[a,b]$\; 


$S(v^c_{0,\pm\delta})\ge r_1(v^c_0)=p+1=S(u_{\pm\delta})+1$. 


Значит, этот случай тоже невозможен. 


 


Таким образом, (38) может быть только в случае (39) при $v^c_0(a)=0$ 


и $v^c_0(b)=0$. Но тогда 


$S(v^c_0)\le r_1(v^c_0,(a,b))=r_1(v^c_0)-2\le p-1$. 


 


Итак, выполняется либо (25), либо (26). 


 


В случае (27) в силу (29), непрерывности $v^c_0(x)$ в $\cal H$ и выбора 


$\varepsilon_0$ 


при достаточно малом $\varepsilon$\; 


$p\ge S(v^{c'}_0,[a+\varepsilon_0,b])\ge S(v^c_0,[a+\varepsilon_0,b])= 


S(v^c_0)=S(u)=p$, 


значит,


\begin{equation} 


S(v^{c'}_0)=p 


\end{equation} 


и $v^{c'}_0(x)$ не имеет узлов в $(a,a+\varepsilon_0)$, 


так что в силу (24) и непрерывности $v^c_0(x)$ в $\cal H$ 


\begin{equation} 


\sign_1v^{c'}_0=\sign_1v^c_0. 


\end{equation} 


Но в силу леммы 1 [12] или теоремы 1 [13] при условии (42) 


$\sign_1v^{c'}_0=\sign_1u$, 


откуда в силу (43) и следует (28). 


\end{pf} 


\setcounter{lem}{1} 


 


\begin{lem} Если $m=1$, $y(\cdot)\in C[a,b]$, 


\begin{equation} 


y(a_1)=0,\quad v(\cdot)=y(\cdot)w(\cdot)\in 


C^{n-r}[a,b] 


\end{equation} 


и выполняются условия $(5)$, $(11)$, то при $a<a_1<b$ 


\begin{equation} 


r_1(y)\le S(L^c_{n-r}v,(a,a_1))+ 


S(L^c_{n-r}v,(a_1,b))+1\le r_1 


(L^c_{n-r}v,(a,b)), 


\end{equation} 


а при $a_1=a$ $(a_1=b)$ 


\begin{equation} 


r_1(y)\le S(L^c_{n-r}v)+1\le r_1(L^c_{n-r}v,[a,b)) 


\ (\le r_1(L^c_{n-r}v,(a,b])). 


\end{equation} 


 


Если при этом $y=\Gamma_c u$, $u(\cdot)\in C_*[a,b]$ и все узлы $u(x)$ 


отличны от $c_j$, $j=\ov{1,r}$, то при $a<a_1<b$ 


\begin{equation} 


S(L^c_{n-r}v,(a,a_1))+S(L^c_{n-r}v,(a_1,b))+1\le S(u), 


\end{equation} 


а при $a_1=a$ или $a_1=b$ 


\begin{equation} 


S(L^c_{n-r}v)+1\le S(u). 


\end{equation} 


\end{lem} 


 


\begin{pf} Так как в силу (17) и обобщенной 


теоремы Ролля (см., напр., лемму 1 [19]) между двумя соседними нулями 


$L^c_{k-1}v$ лежит хотя бы один узел $L^c_k v$, $k=\ov{1,n-r}$, а в 


силу (44) 


$(L^c_k v)(a_1)=0$, $k=\ov 


{0,n-r}$, то при $a<a_1<b$ 


\begin{gather*}  


r_1(y)=r_1(v)\le S(L^c_1v,(a,a_1))+ 


S(L^c_1v,(a_1,b))+1\le r_1(L^c_1v,(a,b)),\\


r_1(L^c_{k-1}v,(a,b))\le S(L^c_k v,(a,a_1))+S(L^c_k v,(a_1,b))+1\le  


r_1(L^c_k v,(a,b)),\quad k=\ov{1,n-r}, 


\end{gather*}  


откуда следует (45). А при $a_1=a$ ($a_1=b$) 


\begin{gather*} 


r_1(y)=r_1(v)\le S(L^c_1v)+1\le r_1(L^c_1v,[a,b)),\\


r_1(L^c_{k-1}v,[a,b))\le S(L^c_k v)+1\le 


r_1(L^c_k v,[a,b)),\quad k=\ov{1,n-r},\\ 


% 


\Big(r_1(y)=r_1(v)\le S(L^c_1v)+1\le r_1 


(L^c_1v,(a,b]),\\


r_1(L^c_{k-1}v,(a,b])\le S(L^c_k v)+1\le r_1(L^c_k v,(a,b]),\quad 


k=\ov{1,n-r}\Big), 


\end{gather*} 


откуда следует (46). 


 


При $y=\Gamma_c u$, $u(\cdot)\in C_*[a,b]$ и отличии всех узлов $u(x)$ 


от $c_j$, $j=\ov{1,r}$, 


(47) и (48) следуют из (18)--(21), (45), (46) и леммы 1. 


\end{pf} 


 


\begin{notenonum} Если $y(\cdot)\in C[a,b]$, 


$v(\cdot)=y(\cdot)w(\cdot)\in C^{n-r}[a,b]$, $y(a_1)\ne 0$, $m=1$, то 


$r_{n-r}(v)\ge n-r+r_1(y)$ 


и, повторяя рассуждения из 


доказательства леммы 8 [10], ввиду (17) получим 


$r_1(y)\le S(L^c_{n-r}v)$. 


\end{notenonum} 


 


\begin{lem} Если все узлы $u(x)$ отличны от $c_j$, $j=\ov{1,r}$, 


то при условиях $(5)$, $(11)$, $(24)$ в случае $m\ge 1$ 


\begin{equation} 


r_1(\Gamma_c u)\le S(u)=p, 


\end{equation} 


а в случае $m=0$ либо выполняется $(49)$, либо 


\begin{equation} 


r_1(\Gamma_c u,(a,b))<p. 


\end{equation} 


 


Если при этом 


\begin{equation} 


S(\Gamma_c u)=S(u), 


\end{equation} 


то 


\begin{equation} 


\sign_1\Gamma_c u=\sign_1 u. 


\end{equation} 


\end{lem} 


 


\begin{pf} При $m\ge 2$ для функций (19)--(21) 


в силу леммы 8 [10] и леммы 1 


\begin{equation} 


S(y_c)=S(\Gamma_c u)\le r_1(\Gamma_c u)=r_1(y_c)\le  


S(L^c_{n-r}v_c)=S(v^c_0)\le 


r_1(v^c_0,(a,b))\le S(u)=p, 


\end{equation} 


т.\,е.~выполняется (49). При $m=1$ в случае $y_c(a_1)\ne 0$ 


(53) и (49) остаются 


верны в силу замечания к лемме 2, а в случае $y_c(a_1)=0$ в силу лемм 


1, 2 


\begin{equation} 


S(y_c)=S(\Gamma_c u)\le r_1(y_c)=r_1(\Gamma_c u)\le 


\max\{r_1(L^c_{n-r}v_c,(a,b]), 


r_1(L^c_{n-r}v_c,[a,b))\}\le S(u)=p, 


\end{equation} 


т.\,е.~(49) тоже выполняется. 


 


При $m\ge 1$ и условии (51) в силу (53) или (54) 


\begin{equation} 


S(y_c)=r_1(y_c)=S(u)=p, 


\end{equation} 


так что $y_c(a)\ne 0$, $y_c(b)\ne 0$, причем для $m\ge 2$ или $m=1$, 


$y_c(a_1)\ne 0$ 


\begin{equation} 


S(y_c)=r_1(y_c)=S(v^c_0)=r_1(v^c_0,(a,b))=S(u)=p. 


\end{equation} 


Если же $m=1$, $y_c(a_1)=0$, то $a<a_1<b$ и $x=a_1$ --- узловой нуль 


$y_c(x)$. 


Поэтому в силу (45), (55) и леммы 1 


\begin{equation} 


S(y_c)=r_1(y_c)=S(L^c_{n-r}v_c,(a,a_1))+ 


S(L^c_{n-r}v_c,(a_1,b))+1= 


r_1(L^c_{n-r}v_c,(a,b))=S(u)=p. 


\end{equation} 


Кроме того, в силу (20)\; $r_{n-r}(v_c,a_1)=n-r$, $v^c_0(a_1)=0$. Если 


$x=a_1$ --- 


пучность $v^c_0(x)$, то в силу (21), (35), (37), (41), (57) при 


достаточно 


малом $\delta$ хотя бы для одной из функций $v^c_{0,\pm\delta}(x)$ 


\begin{equation} 


S(v^c_{0,\pm\delta})\ge p+1>S(u)=S(u_{\pm\delta}), 


\end{equation} 


где все узлы $u_{\pm\delta}(x)$ отличны от $c_j$, $j=\ov{1,r}$.\; (58) 


противоречит лемме 1. Значит, $x=a_1$ может быть только узлом 


$v^c_0(x)$. Но тогда в силу (57) 


снова выполняется (56). 


 


При $m\ge 1$ дальнейшее доказательство дословно повторяет 


доказательство 


леммы 6 [13] и леммы 9 [10]. 


 


Наконец, при $m=0$, т.\,е.~$r=n$, имеем 


$w(x)\equiv1$, $L^c_{n-r}=I$, $v_c(x)\equiv y_c(x)\equiv v^c_0(x)\equiv 


(\Gamma_c u)(x)$ 


и (49) ((50)) следует из (25) ((26)), а (51) 


равносильно (27), так что (52) следует из (28). 


\end{pf} 


 


\begin{lem} При условиях $(5)$, $(11)$ для функций $(19)$--$(21)$ 


\begin{equation} 


S(y_c)\le S(u),\quad S(v^c_0)\le S(u), 


\end{equation} 


в случае же $m=1$, $y_c(a_1)=0$ при $a<a_1<b$ 


\begin{equation} 


S(v^c_0,(a,a_1))+S(v^c_0,(a_1,b))+1\le S(u), 


\end{equation} 


а при $a_1=a$ или $a_1=b$ 


\begin{equation} 


S(v^c_0)+1\le S(u) 


\end{equation} 


$\forall u\in C[a,b]$ и $u\in\delta[a,b]$. 


\end{lem} 


 


\begin{pf} В силу (10), (18)--(21), (33) 


\begin{multline} 


v^c_0(\cdot)=(L^c_{n-r}v_c)(\cdot)=L^c_{n-r}\bigg( 


w(\cdot)\int^b_a\Gamma_c(\cdot,s)u(s)ds\bigg)= \\ 


=\int^b_a L^c_{n-r}\cal G_c(\cdot,s) 


\sigma_c(s)u(s)ds=\int^b_a \cal G_{r c}(\cdot,s) 


\sigma_c(s)u(s)ds. 


\end{multline} 


Ввиду (18), (19), (62) и непрерывности и ограниченности $\Gamma_c(x,s)$ 


и $\cal G_{r c}(x,s)$ в $Q$ любая функция $u(\cdot)\in C[a,b]$ или 


$u(\cdot)\in\delta[a,b]$ может 


быть аппроксимирована такими функциями $u_n(\cdot)\in C_*[a,b]$ с 


отличными 


от $c_j$, $j=\ov{1,r}$, узлами, что 


\begin{equation} 


S(u_n)=S(u)\ \; \forall n\in N,


\end{equation} 


и соответствующие функции $y_{c n}(x)$ ($v^c_{0n}(x)$) сходятся к 


$y_c(x)$ ($v^c_0(x)$) 


равномерно, так что $\exists n_0:S(y_{c n})\ge S(y_c)$, 


$S(v^c_{0n})\ge{}S(v^c_0)$ $\forall n>n_0$. 


Отсюда в силу (63), лемм 1--3 и неравенства 


$S(y,(\alpha,\beta))\le r_1(y,(\alpha,\beta))$ $\forall y\in 


C[\alpha,\beta]$ 


и следуют (59)--(61). 


\end{pf} 


 


\begin{lem} При условиях $(5)$, $(11)$ 


$$ 


r_1(\Gamma_c u)\le S(u)\ \; \forall u\in C[a,b]\ \; 


\text{и} 


\ \; \forall u\in\delta[a,b]. 


$$ 


\end{lem} 


 


\begin{pf} Если $m\ge 2$, то в силу (18)--(21), 


леммы 8 [10] и леммы 4 


\begin{equation} 


r_1(\Gamma_c u)=r_1(y_c)\le S(L^c_{n-r}v_c)= 


S(v^c_0)\le S(u). 


\end{equation} 


При $m=1$, $y_c(a_1)\ne 0$ (64) сохраняется в силу замечания к лемме 2. 


 


Если $m=1$, $y_c(a_1)=0$, то в силу (21) и лемм 2, 4 при $a<a_1<b$ 


$$ 


r_1(\Gamma_c u)=r_1(y_c)\le S(L^c_{n-r}v_c,(a,a_1))+ 


S(L^c_{n-r}v_c,(a_1,b))+1\le S(u), 


$$ 


а при $a_1=a$ или $a_1=b$\; $r_1(\Gamma_c u)=r_1(y_c)\le 


S(L^c_{n-r}v_c)+1\le S(u)$ $\forall u\in C[a,b]$ и $\forall 


u\in\delta[a,b]$. 


\end{pf} 


 


{\bf 4.} Пусть 


\begin{equation} 


\Gamma_c\binom{x_1\dotsb x_k}{s_1\dotsb s_k}= 


\det(\Gamma_c(x_i,s_j))^k_{i,j=1}. 


\end{equation} 


 


Ввиду справедливости леммы 5 далее можно дословно повторить 


рассуждения работ [13], [14] и таким образом доказать следующие 


утверждения. 


 


\begin{lem} При условиях $(5)$, $(11)$ для функции $(10)$ 


$$ 


\Gamma_c\binom{s_1\dotsb s_k}{s_1\dotsb s_k}\ne 0 


\ \;\forall s_i,\ \; i=\ov{1,k},\ \; a<s_1<\dotsb<s_k<b, 


\ \; k=1,2,\dotsc 


$$ 


\end{lem} 


 


\begin{lem} При условиях $(5)$, $(11)$ функция $(10)$ есть 


знакорегулярное 


ядро, т.\,е.~знак каждого ненулевого из определителей $(65)$ 


зависит только от $k$ $\forall x_i$, $s_i$, $i=\ov{1,k}$, $a\le 


x_1<\dotsb<x_k\le b$, $a<s_1<\dotsb<s_k<b$ $\forall k$, $k=1,2,\dotsc$ 


\end{lem} 


 


\begin{lem} При условиях $(5)$, $(11)$ для функции $(10)$  все 


определители $(65)$ 


неотрицательны $\forall x_i$, $s_i$, $i=\ov{1,k}$, $a\le 


x_1<\dotsb<x_k\le b$, $a<s_1<\dotsb<s_k<b$ $\forall k$, $k=1,2,\dotsc$ 


\end{lem} 


 


Леммы 6, 8 в соединении с (12) означают, что при условиях (5), 


(11) $\Gamma_c(x,s)$ есть ядро Келлога на $[\Omega]$ [20], где 


$[\Omega]$ --- ``покомпонентное замыкание'' [20] множества 


$\Omega=\bigcup\limits^r_{j=0}(c_j,c_{j+1})$. 


Можно показать также, что $\Gamma_c(x,s)$ --- осцилляционное ядро по 


терминологии [21]. 


 


\begin{thmnonum} Пусть $q(x)$ суммируема на $[a,b]$, а $\wt 


q(x)=q(x)\sigma_c(x)w(x)>0$ 


почти всюду в $[a,b]$. При $m\ge 1$ и условиях 


$(5)$, $(11)$ спектр задачи $(1)$--$(3)$ осцилляционный, причем для 


собственной 


функции $y_i(x)$ этой задачи, соответствующей собственному значению 


$\lambda_i$, 


все нули функции $\psi_i(x)=y_i(x)/w(x)$ в $[a,b]$ лежат в $(a,b)$ 


$\forall i$, $i=0,1,\dotsc;$ 


любая нетривиальная линейная комбинация 


$\psi_{p k}(x)=\sum\limits^k_{i=p}\alpha_i\psi_i(x)$ 


имеет в $[a,b]$ не менее $p$ узлов и не более $k$ нулей, 


считая пучность 


за два нуля, а при $m\ge 2$ все нули $\psi_i(x)$ простые $\forall i$, 


$i=0,1,\dotsc$, и 


$r_l(\psi_{p k})\le k$, где 


$l=\min\limits_{i=\ov{1,m}}(n-r-k_i)$. 


\end{thmnonum} 


 


{\bf 5.} В случае, когда вместо (5) для некоторой перестановки 


$(j_1,j_2,\dotsc,j_r)$ множества $\{1,2,\dotsc,r\}$, удовлетворяющей 


условиям (8), выполняются 


неравенства (9), положим $\sigma_c(s)=(-1)^{r-j}$ $\forall 


s\in(c_j,c_{j+1})$, если в перестановке 


$(j_1,j_2,\dotsc,j_r)$\; $j$ предшествует $j+1$, и 


$\sigma_c(s)=(-1)^{r-j-1}$ $\forall s\in(c_j,c_{j+1})$ 


в противном случае. При таком определении $\sigma_c(s)$ и условиях (9), 


(11)


в силу теоремы [16] для функции (10) сохраняются неравенства (12). 


Этот факт наводит на предположение, что при замене (5) неравенствами 


(8), (9) предыдущие результаты данной работы сохраняются. При 


некоторых специальных расположениях точек $c_j$, $j=\ov{1,r}$, 


например, при 


\begin{equation} 


c_j=a,\ \; j=\ov{1,q},\ \; c_i=b,\ \; i=\ov{q+1,r}, 


\ \; 0<q<r, 


\end{equation} 


может быть, это и так. По крайней мере, при $r\le 5$ в случае (66) это 


предположение оправдывается. Однако при произвольном расположении 


точек $c_j\in[a,b]$, $j=\ov{1,r}$, и условии (9) вместо (5) с 


$(j_1,j_2,\dotsc,j_r)\ne(1,2,\dotsc,r)$ 


леммы 4, 5, 7, вообще говоря, неверны, о чем свидетельствует 


следующий пример. 


 


Для краевой задачи 


\begin{equation} 


y'''-y''=f(x),\quad y''(0)=0,\quad 


y''(1)-y'(1)=0,\quad y(2)=0 


\end{equation} 


оператор $L=\frac{d^3}{dx^3}-\frac{d^2}{dx^2}$ допускает разложения 


$L=\Big(\frac{d}{dx}-I\Big)L^{1c}_{n-1}=\frac{d}{dx}L^{2c}_{n-1}$, 


$L^{1c}_{n-1}=\frac{d}{dx}L^c_{n-r}$, 


$L^{2c}_{n-1}=\Big(\frac{d}{dx}-I\Big)L^c_{n-r}$, 


$L^c_{n-r}=L^c_1=\frac{d}{dx}$, 


где $n=3$, $r=2$, $c_1=0$, $c_2=1$, так что $\alpha^c_1(x)\equiv1$, 


$\alpha^c_2(x)\equiv0$, $\alpha^c_{12}(x)\equiv0$, 


$\alpha^c_{21}(x)\equiv1$. 


При этом $w(x)=x-2$, $l_{1c_1}y=y''(c_1)=(L^{1c}_{n-1})(c_1)$, 


$l_{2c_2}y=y''(c_2)-y'(c_2)=(L^{2c}_{n-1}y)(c_2)$, $L^{1c}_{n-1}$ и 


$\alpha^c_1(x)$ не зависят от $c_2$ и от $l_{2c_2}$, $L^{2c}_{n-1}$ и 


$\alpha^c_2(x)$ 


не зависят от $c_1$ и от $l_{1c_1}$, $L^c_{n-r}=\frac{d}{dx}$ 


не зависит от порядка 


индексов 1, 2 и $L^c_{n-r}=L^c_1\in T^1_1(a,b)$ $\forall a,b$, $a<b$, а 


$\alpha^c_{12}(x)-\alpha^c_{21}(x)\equiv-1\ne0$ $\forall x$, 


так что в силу теоремы 1 [18] $L\in\wt T^0_3(\cal B;0,2;c_1,c_1)$ 


$\forall c_1,c_2$, $0\le c_1\le c_2\le 1$, 


т.\,е.~$L\in T^0_{31}(\cal B;0,2;0,1)$, где 


$\cal B= 


\left(\begin{smallmatrix}


1 & \phantom{-}0 \\ 1 & -1\end{smallmatrix}\right)$. 


Кроме того, 


$\cal B 


\left(\begin{smallmatrix} 2 & 1 \\ 1 & 2\end{smallmatrix}\right)= 


\left|\begin{smallmatrix}


1 & -1 \\ 1 & \phantom{-}0\end{smallmatrix}\right| 


=1>0$, 


значит, условия (8), (9), (11) выполняются, 


причем можно положить $\sigma_c(s)\equiv1$ в $[0,2]$. 


 


Однако для $u(x)=180x^2-564x+204$ функция $v_c(x)=y_c(x)w(x)$, 


где $y_c(x)=(\Gamma_c u)(x)=\int\limits^2_0\Gamma_c(x,s)u(s)ds$, есть 


решение задачи (67) при 


$f(x)\equiv\sigma(x)u(x)$, т.\,е.~$v_c(x)=-15x^4+34x^3-18x+4$, так что 


$y_c(x)=-15x^3+4x^2+8x-2$, 


а $v^c_0(x)=L^c_{n-r}v_c=\frac{d}{dx}v_c=-60x^3+102x^2-18$. 


Таким образом, $u(x)$ имеет в $[0,2]$ единственный, причем простой, 


нуль 


$x\approx0.417$, тогда как $v^c_0(x)$ имеет в $[0,2]$ простые нули 


$x_1=0.5$, $x_2=0.6+\sqrt{0.96}$, а $y_c(0)<0$, $y_c(0.5)>0$, 


$y_c(1)<0$, так что $r_1(\Gamma_c u)\ge S(\Gamma_c u)=S(y_c)\ge 


2>S(u)=1$ и $S(v^c_0)=2>S(u)=1$. 


Так как теорема 2 [3] может быть доказана для кусочно-непрерывных 


ядер типа (10), то последние неравенства означают, что ядро 


$\Gamma_c(x,s)$ для задачи 


(67) не является знакорегулярным. Итак, для задачи (67) заключения 


лемм 4, 5, 7 не имеют места. 


 


\begin{thebibliography}{99} 


 


\bibitem{1} 


Покорный Ю.В., Лазарев К.П. {\em Некоторые осцилляционные теоремы 


для многоточечных задач} // Дифференц. уравнения. -- 1987. -- 


Т.\,23. -- \No\,4. -- С.\,658--670. 


\bibitem{2} 


Гантмахер Ф.Р., Крейн М.Г. {\em Осцилляционные матрицы и ядра и малые 


колебания механических систем}. -- 2-е изд. -- М.-Л.: Гостехиздат, 


1950. -- 359 с. 


\bibitem{3} 


Левин А.Ю., Степанов Г.Д. {\em Одномерные краевые задачи с операторами, 


не понижающими числа перемен знака}. I // Сиб. матем. журн. -- 1976. 


-- Т.\,17. -- \No\,3. -- С.\,606--626. 


\bibitem{4} 


Левин А.Ю., Степанов Г.Д. {\em Одномерные краевые задачи с операторами, 


не понижающими числа перемен знака}. II // Сиб. матем. журн. -- 1976. 


-- Т.\,17. -- \No\,4. -- С.\,813--830. 


\bibitem{5} 


Покорный Ю.В. {\em О функции Грина многоточечной краевой задачи} // 


Дифференц. уравнения. -- 1978. -- Т.\,14. -- \No\,4. -- С.\,760--761. 


\bibitem{6} 


Дерр В.Я. {\em О спектре некоторых многоточечных задач}. -- Удм. гос. 


ун-т, Устинов. мех. ин-т. -- Устинов, 1987. -- 46 с. -- 


Деп. в ВИНИТИ 23.01.87, \No\,533-В87. 


\bibitem{7} 


Покорный Ю.В., Шурупова И.Ю. {\em Об осцилляционных свойствах спектра 


краевой задачи с функцией Грина, меняющей знак} // Укр. матем. журн. -- 


1989. -- Т.\,41. -- \No\,11. -- С.\,1521--1526. 


\bibitem{8} 


Степанов Г.Д. {\em Многоточечные краевые задачи с функциями Грина, 


приводимыми к знакорегулярному виду}. -- Ярослав. ун-т. -- 


Ярославль, 1988. -- 


25 с. -- Деп. в ВИНИТИ 20.04.88, \No\,3044-В88. 


\bibitem{9} 


Тептин А.Л. {\em Об осцилляционных свойствах спектра одной краевой 


задачи} // Функц.-дифференц. уравнения. -- Пермь, 1987. -- С.\,112--124. 


\bibitem{10} 


Тептин А.Л. {\em Об осцилляционности ядра, связанного с функцией 


Грина одной многоточечной задачи}. -- 


Ред. журн. ``Дифференц. уравнения''. -- Минск, 


1989. -- 30 с. -- Деп. в ВИНИТИ 11.08.89, \No\,5417-В89. 


\bibitem{11} 


Дерр В.Я. {\em О спектре задачи с комбинированными краевыми условиями} 


// Функц.-дифференц. уравнения. -- Пермь, 1987. -- С.\,105--112. 


\bibitem{12} 


Тептин А.Л. {\em Об осцилляционности $($с точностью до множителя$)$ 


функции Грина одной многоточечной задачи} // Функц.-дифференц. 


уравнения. -- Пермь, 1992. -- С.\,159--172. 


\bibitem{13} 


Тептин А.Л. {\em Об осцилляционности спектра краевой задачи с функцией 


Грина, меняющей знак по обоим аргументам}. -- Ижевск. механ. ин-т. -- 


Ижевск, 1993. -- 49 с. -- Деп. в ВИНИТИ 12.03.93, \No\,596-В93. 


\bibitem{14} 


Тептин А.Л. {\em Об осцилляционности спектра одной многоточечной 


краевой задачи} // Дифференц. уравнения. -- 


1995. -- Т.\,31. -- \No\,8. -- С.\,1370--1380. 


\bibitem{15} 


Тептин А.Л. {\em O многоточечной краевой задаче, функция Грина которой 


меняет знак в ``шахматном'' порядке} // Дифференц. 


уравнения. -- 1984. -- 


Т.20. -- \No\,11. -- С.\,1910--1914. 


\bibitem{16} 


Тептин А.Л. {\em Новый признак знакопостоянства функции Грина} // 


Дифференц. уравнения. -- 1988. -- Т.\,24. -- \No\,6. -- 


С.\,1066--1069. 


\bibitem{17} 


Левин А.Ю. {\em Неосцилляция решений уравнения 


$x^{(n)}+p_1(t)x^{(n-1)}+\dotsb+p_n(t)x=0$} 


// УМН. -- 1969. -- Т.\,24. -- Вып.\,2. -- C.\,43--96. 


\bibitem{18} 


Тептин А.Л. {\em О неосцилляции решений и знаке функции Грина} // 


Дифференц. уравнения. -- 1984. -- Т.\,20. -- \No\,6. -- С.\,995--1005. 


\bibitem{19} 


Чичкин Е.С. {\em Теорема о дифференциальном неравенстве для 


многоточечных краевых задач} // Изв. вузов. Математика. -- 1962. 


-- \No\,2. -- С.\,170--179. 


\bibitem{20} 


Покорный Ю.В., Боровских А.В. {\em О теореме Келлога для разрывных 


функций Грина} // Матем. заметки. -- 1993. -- Т.\,53. -- \No\,1. -- 


С.\,151--153. 


\bibitem{21} 


Боровских А.В. {\em О спектре интегральных операторов с разрывными 


ядрами}. -- Воронеж. ун-т. -- Воронеж, 1989. -- 27 с. -- 


Деп. в ВИНИТИ 15.03.89, \No\,1684-В89. 


 


\end{thebibliography} 


 


\vskip 0.5cm 


 


\post{\it Ижевский государственный}{\it Поступила} 


\post{\it технический университет}{11.07.1996} 


 


\label{end} 


\end{document} 


 


